\bookchapter{Lampiran G: Indeks Panduan Cepat Seluruh Buku}

\chapterintro{Diurutkan berdasarkan pinyin agar kamu dengan cepat menemukan bab yang dituju.}

\booksubsection{A-B}

\begin{itemize}
  \item Kunci BL Bab 24
  \item ADB Bab 20
  \item Instalasi dan Pengelolaan App Bab 4
  \item Asisten AI Bab 22
  \item Ponsel AI Bab 26
  \item Penipuan Wajah Palsu AI Bab 32
  \item Pemrograman AI Bab 37-47
  \item Lima Prinsip Pemrograman AI Bab 38
  \item API Key Bab 42
  \item Agent (Agen Cerdas) Bab 37
  \item CLAUDE.md Bab 39
  \item Pengembangan Android Bab 48-58
  \item Android Studio Bab 49
  \item Siklus Hidup Activity Bab 52
  \item Menerbitkan ke App Store Bab 67
\end{itemize}

\booksubsection{C-D}

\begin{itemize}
  \item Konektivitas Ponsel-Mobil Bab 34
  \item Baterai dan Daya Tahan Bab 15
  \item Pengelolaan Kata Sandi Bab 14
  \item Passkey Bab 14
  \item Mode Pengembang Bab 20
  \item Anak-anak dan Ponsel Bab 35
  \item Claude Code / Cursor Bab 41
  \item DeepSeek Bab 42
  \item Compose Bab 51
  \item DataStore Bab 53
\end{itemize}

\booksubsection{F-H}

\begin{itemize}
  \item Layar Lipat (Foldable) Lampiran F
  \item Pemecahan Masalah Bab 19
  \item Menerbitkan ke Google Play Bab 56
  \item Pengisi Daya GaN Bab 15
  \item iCloud Bab 16
  \item Rumah Pintar Bab 36
  \item Fotografi Komputasi Tambahan Bab 11
\end{itemize}

\booksubsection{K-R}

\begin{itemize}
  \item Teknik Memotret Bab 11
  \item Pengelolaan Izin Bab 4
  \item Sertifikasi SafetyNet Bab 20
  \item Serangan SIM Swap Bab 32
  \item Akses Root Bab 20 / Bab 24
  \item Termux Tambahan Bab 20
  \item Protokol MCP Bab 46
  \item Dokumen Kebutuhan PRD Bab 40
  \item Prompt Engineering Bab 43
  \item Kotlin Bab 50
  \item Arsitektur MVVM Bab 57 · Bab 68
  \item Database Room Bab 53
  \item Retrofit Bab 53
\end{itemize}

\booksubsection{S-Z}

\begin{itemize}
  \item Warisan Digital Bab 25
  \item Cadangan Data Bab 18
  \item Penyaringan Gangguan Bab 8 · Bab 13
  \item Peringatan Cuaca Bab 13
  \item Peringatan Bencana Alam Bab 13 · Bab 32
  \item Peringatan Gempa Bab 13 · Bab 32
  \item Aplikasi Berguna dan Alat Pembersih Bab 13
  \item Qingzhuo Bab 13
  \item UWB (Ultra-Wideband) Bab 17
  \item Keamanan Jaringan Bab 21
  \item Panduan Membeli Ponsel Bab 28
  \item Layanan Cloud Bab 16
  \item Pengembangan Berbasis Spesifikasi (SDD) Bab 39
  \item Token Bab 42
  \item Vibe Coding Bab 37
  \item Swift Bab 61
  \item SwiftUI Bab 62
  \item SwiftData Bab 64
  \item URLSession Bab 64
  \item Xcode Bab 60
  \item Pengembangan iOS Bab 59-69
\end{itemize}

\booksubsection{Lainnya}

\begin{itemize}
  \item Pengelolaan File Bab 2
  \item Fitur Ramah Lansia Bab 27
  \item Rekomendasi App Perkakas Bab 13
  \item Anti-Gangguan / National Anti-Fraud Center Bab 13 · Bab 27
  \item Remote Komputer (Sunflower/ToDesk) Bab 13
  \item Kantor AI (Trae dll.) Bab 13 · Bab 41
  \item LocalSend / Kirim Cepat Jaringan Lokal Bab 13
  \item KDE Connect / Sinkronisasi Papan Klip Bab 13 · Bab 17
  \item Kata Sandi Dinamis 2FA (Aegis dll.) Bab 13 · Bab 14
  \item TickTick Bab 13
  \item Pemindaian Dokumen Kertas Bab 13
  \item Artificial Desktop (Wallpaper Dinamis) Bab 13
  \item Darurat Ponsel Saat Gempa/Bencana Bab 32
  \item Screencasting dan Mobil Bab 30
  \item Pembuatan Video Pendek Bab 31
  \item P3K Keamanan Bab 32
  \item Git/GitHub Tambahan Bab 23
  \item Kontrol Versi Git Bab 45
  \item Context Engineering Bab 39
  \item Requirements Engineering Bab 40
  \item App Pembukuan Bab 54 · Bab 65
  \item App Cuaca Bab 55 · Bab 66
  \item Widget WidgetKit Bab 68
\end{itemize}
